% Abschnitt 3: Beispielrechnung (mit Tabelle)
\section{Beispielrechnung}
\label{sec:beispiel}

Tabelle~\ref{tab:beispiel} listet die Ressourcen einer Beispielseite und ihre
Übertragungszeit $S/B$ bei einer Datenrate von
$B = \SI{10}{\mega\bit\per\second}$ auf (\SI{1}{\kilo\byte} entspricht
\SI{8}{\kilo\bit}). Die Bilder verursachen fast drei Viertel der
Übertragungszeit und sind damit der wirksamste Ansatzpunkt.

\begin{table}[htb]
  \centering
  \caption{Ressourcen einer Beispielseite und Übertragungszeit bei
           \SI{10}{\mega\bit\per\second}.}
  \label{tab:beispiel}
  \begin{tabular}{lrr}
    \toprule
    Ressource & Datenmenge $S$ & Übertragungszeit $S/B$ \\
    \midrule
    HTML-Dokument     & \SI{25}{\kilo\byte}   & \SI{20}{\milli\second}   \\
    Stylesheets (CSS) & \SI{60}{\kilo\byte}   & \SI{48}{\milli\second}   \\
    JavaScript        & \SI{350}{\kilo\byte}  & \SI{280}{\milli\second}  \\
    Bilder            & \SI{1200}{\kilo\byte} & \SI{960}{\milli\second}  \\
    \midrule
    Summe             & \SI{1635}{\kilo\byte} & \SI{1308}{\milli\second} \\
    \bottomrule
  \end{tabular}
\end{table}

Mit den angenommenen Werten $t_\mathrm{DNS} = \SI{20}{\milli\second}$,
$t_\mathrm{V} = \SI{60}{\milli\second}$ und
$t_\mathrm{S} = \SI{100}{\milli\second}$ liefert
Gleichung~\eqref{eq:ladezeit} eine Ladezeit von
$t_\mathrm{L} = \SI{1488}{\milli\second}$, also rund \SI{1,5}{\second}.
Halbiert man die Datenmenge der Bilder, sinkt sie auf etwa \SI{1,0}{\second}.
